%% r2_blochkato_paper.tex
%% "Bloch-Kato Tamagawa conjecture for M(4.5.b.a) at interior critical integers,
%%  with a Damerell-Tamagawa ratio invariant"
%% Sub-agent M70 (Sonnet 4.6), 2026-05-06
%% Status: DRAFT (5 [TBD: prove] markers, ~10-15% probability of formal contribution)
%% Target: Research in Number Theory (primary) / Journal of Number Theory (backup)

\documentclass[12pt]{amsart}

\usepackage{amsmath,amssymb,amsthm,mathrsfs}
\usepackage[colorlinks=true,linkcolor=blue,citecolor=blue,urlcolor=blue]{hyperref}
\usepackage{microtype}
\usepackage{enumitem}

%% Theorem environments
\newtheorem{conjecture}{Conjecture}[section]
\newtheorem{theorem}[conjecture]{Theorem}
\newtheorem{proposition}[conjecture]{Proposition}
\newtheorem{remark}[conjecture]{Remark}
\theoremstyle{definition}
\newtheorem{definition}[conjecture]{Definition}
\newtheorem{example}[conjecture]{Example}

%% Shorthand
\newcommand{\Z}{\mathbb{Z}}
\newcommand{\Q}{\mathbb{Q}}
\newcommand{\R}{\mathbb{R}}
\newcommand{\C}{\mathbb{C}}
\newcommand{\Zp}{\mathbb{Z}_p}
\newcommand{\Qp}{\mathbb{Q}_p}
\newcommand{\A}{\mathbb{A}}
\newcommand{\Gal}{\mathrm{Gal}}
\newcommand{\Hom}{\mathrm{Hom}}
\newcommand{\End}{\mathrm{End}}
\newcommand{\Sel}{\mathrm{Sel}}
\newcommand{\Sha}{\Sha}
\newcommand{\Tam}{\mathrm{Tam}}
\newcommand{\ord}{\mathrm{ord}}
\newcommand{\lcm}{\mathrm{lcm}}
\newcommand{\Spec}{\mathrm{Spec}}
\newcommand{\OmegaK}{\Omega_K}
\newcommand{\Omegapm}{\Omega^\pm}
\newcommand{\tbd}[1]{\textbf{[TBD: prove]} \emph{#1}}
\newcommand{\tbdtag}[2]{\textbf{[TBD-R2-#1: prove]} \emph{#2}}

\begin{document}

\title{Bloch-Kato Tamagawa conjecture for $M(f)$ at interior critical integers,\\
with a Damerell-Tamagawa ratio invariant}

\author{K.~Remondiere}
\address{Independent researcher, affiliated with crossed-cosmos ECI project}
\email{kevin.remondiere@gmail.com}

\date{2026-05-06 (draft)}

\begin{abstract}
Let $f = \mathbf{4.5.b.a}$ (LMFDB) be the unique weight-$5$, level-$4$, CM newform
with CM field $K = \Q(i)$ and Steinberg edge eigenvalue $a_2 = -4$.
We formulate three conjectures (R-2.1, R-2.2, R-2.3)
on the Bloch-Kato Tamagawa number formula for the Scholl motive
$M(f)$ at the interior critical integers $m \in \{1,2,3,4\}$.
The central observation is an $\Omega$-\emph{independent} ratio
\[
  \pi \cdot L(f,1) \,/\, L(f,2) \;=\; \tfrac{6}{5},
\]
verified to 80-digit precision in PARI/GP and cross-checked in SageMath.
Conjecture~R-2.2 proposes that this ratio factors as a product of
local Tamagawa factors at $p = 2$, Bloch-Kato Selmer quotients, and archimedean
Gamma-ratio factors, yielding an intrinsic Diophantine fingerprint of $M(f)$.
We document the complete obstruction: every published Tamagawa Number Conjecture
framework from 2024-2026 (Büyükboduk-Neamti, Castella, BBL, Sano, Yin,
Diamond-Flach-Guo) is \emph{blocked} by at least one hypothesis violation for the
regime $(p=2 \text{ ramified},\ k=5 \text{ odd},\ K=\Q(i),\ N=4)$.
Five \textbf{[TBD: prove]} markers are stated explicitly.
We estimate a $\sim 10$-$15\%$ probability that a formal Bloch-Kato Tamagawa
contribution is achievable without a new $p=2$-ramified anticyclotomic Iwasawa
framework.
\end{abstract}

\maketitle

\tableofcontents

%% ============================================================
\section{Introduction}
\label{sec:intro}
%% ============================================================

The Bloch-Kato conjecture \cite{BK90}, in its Tamagawa number formulation
\cite{FPR94}, predicts that for a pure motive $M$ over $\Q$ with
good or semistable reduction, the leading term of the $L$-function $L(M,s)$
at a critical integer $m$ is related, modulo periods and powers of $2\pi i$,
to an arithmetic ratio
\begin{equation}\label{eq:BK}
  \frac{L^*(M(m),0)}{\Omega^\pm(M(m))}
  \;=\;
  \frac{\#H^0(\Q, V(m)) \cdot \prod_\ell c_\ell(M(m))}{\#H^1_f(\Q, V(m))^{-1} \cdot \#H^2_f(\Q, V(m))^{-1}},
\end{equation}
where $V$ is the $p$-adic realization, $H^1_f$ is the Bloch-Kato Selmer group,
$c_\ell$ are local Tamagawa factors at all primes $\ell$, and $\Omega^\pm$
are Deligne's canonical periods \cite{Del79}.

\medskip
The higher-weight BSD-analogue for CM elliptic modular forms was reopened
by two recent developments in the crossed-cosmos ECI programme:
\begin{enumerate}[label=(\roman*)]
\item \textbf{The M52 Damerell ladder} \cite{Dam70,Dam71}: for $f = 4.5.b.a$,
  the algebraic parts $\alpha_m = L(f,m)/\Omega_K^{4} \in \Q$ for all
  $m \in \{1,2,3,4\}$, verified to 80-digit precision in PARI/GP.
  Among all CM weight-$5$ dim-$1$ newforms tested, this full-$\Q$ rationality
  is \emph{unique} to $4.5.b.a$ (Q($\omega$) competitors yield $\alpha_m \in \Q(\sqrt{3}) \setminus \Q$ for odd $m$).

\item \textbf{The M57 Adelic Katz conjecture} \cite{Katz78,Kriz21}:
  an anticyclotomic 2-adic $L$-function $\mathscr{L}_2^\pm(f)$
  is conjectured to exist in $D(\Gamma, \Z_2)$, interpolating the Damerell
  values with Frobenius-correction factors.
\end{enumerate}

\medskip
The $\Omega$-independent ratio $\pi \cdot L(f,1)/L(f,2) = 6/5$
(Section~\ref{sec:damerell}) is the key numerical input. It is \emph{independent}
of all period choices and is therefore a purely Diophantine quantity.
Conjecture~R-2.2 (Section~\ref{sec:conjectures}) proposes that it factors
via a product of local Tamagawa quotients and Selmer quotients, making the $6/5$
the \emph{central novel content} of this note.

\medskip
\textbf{Honest scope.} We estimate the probability of a formal Bloch-Kato
Tamagawa contribution at $\sim 10$-$15\%$, conditional on
the development of a new $p=2$-ramified anticyclotomic Iwasawa framework.
The complete obstruction is documented in Section~\ref{sec:blocked}.
The five \textbf{[TBD: prove]} markers (Section~\ref{sec:open}) are explicit
and not hidden.

\medskip
\textbf{Acknowledgements.} The numerical computations used PARI/GP \cite{PARI}
and SageMath \cite{Sage}. LMFDB data \cite{LMFDB} provided Hecke eigenvalues,
CM field, and conductor.

%% ============================================================
\section{The motive \texorpdfstring{$M(f)$}{M(f)} of \texorpdfstring{$f = 4.5.b.a$}{f = 4.5.b.a}}
\label{sec:motive}
%% ============================================================

\subsection{LMFDB data for \texorpdfstring{$f = 4.5.b.a$}{f = 4.5.b.a}}

The newform $f$ has the following LMFDB-verified invariants \cite{LMFDB}:
\begin{itemize}
  \item Weight $k = 5$, level $N = 4 = 2^2$, character $\chi_{-4}$ (primitive, conductor $4$).
  \item Hecke eigenvalues: $a_2 = -4$, $a_3 = 0$, $a_5 = -20$,
    $a_7 = 0$, $a_{13} = 0$, $a_{17} = 100$ (all CM vanishing at inert primes).
  \item CM field: $K = \Q(i) = \Q(\sqrt{-1})$, discriminant $d_K = -4$.
  \item Inducing Hecke Grössencharacter $\psi$ of conductor $(1+i)^2 = (2)$ in $\Z[i]$
    with infinity type $(k-1, 0) = (4, 0)$.
  \item Dimension of space $S_5(\Gamma_1(4), \chi_{-4})$: dim $= 1$ (Cohen-Oesterlé \cite{CO77}).
  \item Analytic conductor: $4$; root number $\varepsilon(f) = +1$.
  \item Minimal twist: itself (all level-$36$, $64$, $100$ CM weight-$5$ forms are twists of $4.5.b.a$).
\end{itemize}

The Steinberg edge condition $a_2 = -4 = \varepsilon_2 \cdot 2^{(k-1)/2}$ with $\varepsilon_2 = -1$
is unique among CM weight-$5$ dim-$1$ newforms with $K = \Q(i)$ (the twist $36.5.d.a$
has $a_2 = +4$, i.e., $\varepsilon_2 = +1$; forms of level $64$ and $100$ are non-Steinberg at $2$).
This Steinberg-with-$\varepsilon_2=-1$ structure drives the alternating sign pattern
in the 2-adic valuation monotone ladder (Section~\ref{sec:damerell}).

\subsection{Scholl motive}

By Scholl's theorem \cite{Sch90}, there exists a pure Chow motive
\[
  M(f) \;\in\; \mathrm{Mot}(\Q)_\Q
\]
of rank $2$, weight $w = k-1 = 4$, attached to $f$. Its realizations are:
\begin{itemize}
  \item \emph{Betti:} $M(f)_B = H^1_B(X,\Q)$ for a suitable variety $X/\Q$,
    with Hodge decomposition $M(f)_B \otimes \C = H^{4,0} \oplus H^{0,4}$.
  \item \emph{de Rham:} $M(f)_{dR}$ with Hodge filtration
    $\mathrm{Fil}^4 M(f)_{dR} \neq 0$, $\mathrm{Fil}^5 M(f)_{dR} = 0$.
  \item \emph{$\ell$-adic:} $V_\ell(f) = $ the $\ell$-adic representation attached to $f$
    by Deligne \cite{Del71}, a rank-$2$ $\Q_\ell$-representation of $G_\Q = \Gal(\overline{\Q}/\Q)$.
  \item \emph{$2$-adic:} $V_2(f)$, a rank-$2$ $\Q_2$-representation.
    At $p = 2$: since $2 \mid N$ with $\ord_2(N) = 2$ and $a_2 = -4$,
    the local $(\varphi, N)$-module at $p=2$ is of Steinberg type;
    $V_2(f)|_{G_{\Q_2}}$ is an extension of unramified-twist
    $\Q_2(-1)$ by $\Q_2(k-1)\chi$, where $\chi$ is the unramified quadratic character.
\end{itemize}

\subsection{Interior critical integers and Deligne periods}

The motivic weight $w = 4$ places the central point at $s = w/2 + 1 = 3$.
The critical integers for $L(M(f), s)$ are $s = 1, 2, 3, 4$, i.e.,
$m \in \{1,2,3,4\}$ in the notation $L(M(f)(m), 0) = L(f, m)$.

The Deligne period \cite{Del79} decomposes as
\[
  \Omega^\pm(M(f)(m)) = \Omega_K^{2m} \cdot (2\pi i)^{-(k-1-m)},
\]
where $\Omega_K = \Gamma(1/4)^2/(2\sqrt{2\pi})$ is the lemniscate period
of $K = \Q(i)$, normalized via the Chowla-Selberg formula \cite{CS67}.
Here $\Gamma(1/4)$ is the classical Gamma function value; numerically
$\Omega_K \approx 2.6220575543\ldots$

The algebraic normalization from Shimura \cite{Shim76} and Damerell
\cite{Dam70, Dam71} yields $\alpha_m := L(f,m)/\Omega_K^4 \in \Q$,
verified in Section~\ref{sec:damerell}.

%% ============================================================
\section{Damerell ladder and the \texorpdfstring{$6/5$}{6/5} invariant}
\label{sec:damerell}
%% ============================================================

\subsection{The Damerell ladder (PARI 80-digit, Sage cross-check)}

The \emph{Damerell ladder} consists of the four algebraic special values
$\alpha_m = L(f,m)/\Omega_K^4$, $m \in \{1,2,3,4\}$.

\begin{theorem}[M52 PARI 80-digit verification]
\label{thm:damerell}
For $f = 4.5.b.a$ and $\Omega_K = \Gamma(1/4)^2/(2\sqrt{2\pi})$:
\[
  \alpha_1 = \frac{1}{10}, \quad
  \alpha_2 = \frac{1}{12}, \quad
  \alpha_3 = \frac{1}{24}, \quad
  \alpha_4 = \frac{1}{60}.
\]
Each equality holds to within $10^{-60}$ in PARI/GP at 80-digit precision.
\end{theorem}

\begin{remark}
The PARI computation proceeds via numerical integration of $L(f,m)$
using the built-in \verb|lfun()| function, followed by division by
$\Omega_K^4$ computed from the Chowla-Selberg formula.
A Stage~1 SageMath cross-check via
\begin{verbatim}
  f = Newforms(Gamma1(4), 5, names='a')[0]
  f.lseries()(m)
\end{verbatim}
yields agreement to default precision ($\approx 10^{-16}$),
confirming the PARI values. This cross-check is at default
Sage precision and should not be treated as an independent high-precision verification.
\end{remark}

\subsection{Uniqueness of the full-\texorpdfstring{$\Q$}{Q} ladder}

Among CM weight-$5$ dim-$1$ newforms, the full $\Q$-rationality
$\alpha_m \in \Q$ for all $m \in \{1,2,3,4\}$ is \emph{unique} to $4.5.b.a$.

For $K = \Q(\sqrt{-3})$ (Eisenstein CM), the analogous values satisfy:
$\alpha_m \in \Q(\sqrt{3}) \setminus \Q$ for $m$ odd, because the
Hecke Grössencharacter has order $6$ (units $\{\pm 1, \pm\omega, \pm\omega^2\}$)
and $\mathrm{Im}(\omega) = \sqrt{3}/2$ contaminates odd-$m$ periods.
Specifically, for $27.5.b.a$ ($Q(\omega)$): $\alpha_1 = 27\sqrt{3}/4$ and
$\alpha_3 = \sqrt{3}/4$; for $12.5.c.a$: $\alpha_1 = 2\sqrt{3}$ and
$\alpha_3 = 2\sqrt{3}/9$.

For $K = \Q(i)$: the lemniscate period $\Omega_K = \Gamma(1/4)^2/(2\sqrt{2\pi})$
is $\sqrt{3}$-free, so all four algebraic parts land in $\Q$.
This is a unique property of the class-number-$1$ imaginary quadratic field
with $d_K = -4$.

\subsection{The \texorpdfstring{$\Omega$}{Omega}-independent \texorpdfstring{$6/5$}{6/5} invariant}

\begin{proposition}[M52, $\Omega$-independent]
\label{prop:sixfifths}
The ratio
\[
  \rho := \pi \cdot \frac{L(f,1)}{L(f,2)}
\]
is independent of the period $\Omega_K$ and equals
\[
  \rho = \frac{\alpha_1}{\alpha_2} \cdot \pi
        \;=\; \frac{1/10}{1/12} \cdot \pi \cdot \frac{\text{[archimedean factor]}}{\text{[archimedean factor]}}
        \;=\; \frac{6}{5}.
\]
More precisely, with $L(f,m) = \alpha_m \cdot \Omega_K^4 \cdot (2\pi)^{m-(k-1)} \cdot \mathrm{[Gamma]}$,
one finds $\pi \cdot L(f,1)/L(f,2) = 6/5 \in \Q$,
verified to 80-digit precision in PARI/GP and to default precision ($\sim 10^{-16}$) in SageMath.
\end{proposition}

\begin{remark}
The analogous ratio for $Q(\omega)$ competitors:
$\pi \cdot L(f,1)/L(f,2) = 3\sqrt{3} \notin \Q$ for $27.5.b.a$
and $= 3\sqrt{3}/2 \notin \Q$ for $12.5.c.a$.
The rationality of $\rho = 6/5$ is therefore an $\Omega$-independent diagnostic
criterion separating $4.5.b.a$ from all tested CM weight-$5$ competitors.
\end{remark}

\subsection{F1 Steinberg renormalization and the monotone ladder}

The local factor at $p = 2$ (Steinberg, $\varepsilon_2 = -1$)
requires a Frobenius-correction factor (``F1'' renormalization):
\[
  E_2^\pm(f, m) = (-2^{m-1}) \cdot (1 + 2^{m-3}),
\]
yielding renormalized values $\alpha_m^{F1} = E_2^\pm(f,m) \cdot \alpha_m$.
The 2-adic valuations are
\[
  v_2(\alpha_m^{F1}) = \{-3, -2, 0, +1\} \quad (m = 1, 2, 3, 4).
\]
This \emph{monotone} alternating fingerprint is verified by sympy (M22)
and PARI (M22, F2~v3, F2~v5), and is \emph{unique} to $4.5.b.a$
among level-$4$ CM weight-$5$ forms.

%% ============================================================
\section{Conjectures R-2.1, R-2.2, R-2.3}
\label{sec:conjectures}
%% ============================================================

We now state three interrelated conjectures.
All are \emph{open}; none is claimed as a theorem.

\subsection{Setup}

Let $p = 2$, $T_2(f) \subset V_2(f)$ a $G_\Q$-stable $\Z_2$-lattice,
and for $m \in \{1,2,3,4\}$ write $T_2(f)(m)$ for the $m$-th Tate twist.
Define:
\begin{itemize}
  \item $H^1_f(\Q, V_2(f)(m))$: the Bloch-Kato Selmer group \cite{BK90},
    computed using the Fontaine-Laffaille / Fontaine-Perrin-Riou condition \cite{FPR94}.
  \item $c_2(M(f)(m))$: the local Tamagawa factor at $\ell = 2$
    (the only ramified prime since $N = 4 = 2^2$).
  \item $c_\infty(M(f)(m))$: the archimedean Gamma-factor contribution.
  \item $\Omega^\pm_m = \Omega^\pm(M(f)(m))$: the Deligne period.
  \item $\mathrm{Tam}(M(f)(m)) = c_2(M(f)(m)) \cdot c_\infty(M(f)(m))$:
    total Tamagawa factor (only $\ell = 2$ contributes non-trivially at finite places).
\end{itemize}

\subsection{Conjecture R-2.1}

\begin{conjecture}[R-2.1: BK Tamagawa formula]
\label{conj:R21}
For each $m \in \{1,2,3,4\}$:
\begin{equation}\label{eq:R21}
  \frac{L^*(M(f)(m), 0)}{\Omega^\pm_m}
  \;=\;
  \frac{\# H^1_f(\Q, V_2(f)(m))}{\# H^0(\Q, V_2(f)(m)) \cdot \# H^2_f(\Q, V_2(f)(m))}
  \;\cdot\;
  c_2(M(f)(m)) \cdot c_\infty(M(f)(m)),
\end{equation}
where the left-hand side equals $\alpha_m$ (Theorem~\ref{thm:damerell}).
\end{conjecture}

\begin{remark}
Since $f$ is non-ordinary at $p = 2$ (supersingular-like with $v_2(a_2) = 2$)
and $2 \mid N$, one expects $H^0(\Q, V_2(f)(m)) = 0$ and $H^2_f = H^0(\Q, V_2(f)^*(1-m))^\vee = 0$
for $1 \leq m \leq 4$, by a parity argument. Under this vanishing,
\eqref{eq:R21} simplifies to $\alpha_m = c_2(m) \cdot c_\infty(m) / \#H^1_f(\Q, V_2(f)(m))$.
\end{remark}

\tbdtag{1}{Compute the local Tamagawa factor $c_2(M(f)(m))$
at the ramified Steinberg prime $p = 2$ via Tate duality
and local $\varepsilon$-factors. The Steinberg representation
at $p = 2$ (with $\varepsilon_2 = -1$, $\ord_2(N) = 2$) lies
outside the range of Fontaine-Laffaille theory and requires
direct Tate local cohomology computation.}

\tbdtag{2}{Compute $\#H^1_f(\Q, V_2(f)(m))$ for $m = 1, 2, 3, 4$.
No Heegner-point machinery (Kolyvagin, BDP, Castella)
is available in the $(p=2 \text{ ramified}, K=\Q(i))$ regime
(see Section~\ref{sec:blocked}).
Expected heuristic: $H^1_f$ is finite for CM motives by analogy
with Rubin's theorem \cite{Rub91} for CM elliptic curves,
but this requires a new anticyclotomic descent at $p = 2$.}

\subsection{Conjecture R-2.2}

\begin{conjecture}[R-2.2: Tamagawa-ratio factorization of $6/5$]
\label{conj:R22}
The $\Omega$-independent invariant $\rho = 6/5$ admits a factorization
\begin{equation}\label{eq:R22}
  \frac{6}{5}
  \;=\;
  \frac{\alpha_1}{\alpha_2}
  \;=\;
  \frac{c_2(M(f)(1))}{c_2(M(f)(2))}
  \;\cdot\;
  \frac{\#H^1_f(\Q, V_2(f)(2))}{\#H^1_f(\Q, V_2(f)(1))}
  \;\cdot\;
  \frac{c_\infty(M(f)(1))}{c_\infty(M(f)(2))},
\end{equation}
where the three factors are:
\begin{itemize}
  \item $c_2(1)/c_2(2)$: ratio of local Tamagawa factors at the ramified prime $2$;
  \item $\#H^1_f(2)/\#H^1_f(1)$: inverse ratio of Selmer group orders;
  \item $c_\infty(1)/c_\infty(2)$: archimedean Gamma-ratio from Deligne's period comparison.
\end{itemize}
Heuristically, the numerator ``$6$'' arises from $c_2(1) \cdot \#H^1_f(2)$
and the denominator ``$5$'' from $c_2(2) \cdot \#H^1_f(1)$,
up to the archimedean correction.
\end{conjecture}

\tbdtag{3}{Establish the factorization \eqref{eq:R22} explicitly.
This is the \emph{central novel content} of Conjectures R-2.
Neither factor alone can be computed without resolving TBD-R2-1 and TBD-R2-2.
The identity $6/5 = (\text{Tamagawa-ratio}) \times (\text{Selmer-ratio}) \times (\text{Gamma-ratio})$
is the primary open problem of this note; its verification
would constitute a non-trivial test of the Bloch-Kato formula \eqref{eq:R21}
for a CM motive at a Steinberg-edge ramified prime.}

\subsection{Conjecture R-2.3}

\begin{conjecture}[R-2.3: 2-adic Tamagawa and the monotone ladder]
\label{conj:R23}
For $m \in \{1,2,3,4\}$:
\begin{equation}\label{eq:R23}
  v_2(c_2(M(f)(m))) - v_2(\#H^1_{f,2}(\Q, V_2(f)(m)))
  \;=\;
  v_2(\alpha_m^{F1}) \;\in\; \{-3, -2, 0, +1\},
\end{equation}
where $\alpha_m^{F1} = E_2^\pm(f,m) \cdot \alpha_m$ is the
F1-renormalized Damerell value (Section~\ref{sec:damerell}),
and $H^1_{f,2}$ denotes the $2$-part of the Bloch-Kato Selmer group.
\end{conjecture}

\begin{remark}
The monotone sequence $v_2(\alpha_m^{F1}) = \{-3, -2, 0, +1\}$ for
$m = \{1, 2, 3, 4\}$ was verified unconditionally by sympy (M22)
and numerically by PARI (M22, F2~v3, F2~v5). Conjecture~R-2.3
proposes that this arithmetic sequence has a \emph{motivic explanation}
in terms of the relative sizes of local Tamagawa factors and Selmer groups.
\end{remark}

\tbdtag{4}{Verify Conjecture~R-2.3 assuming TBD-R2-1 and TBD-R2-2.
The predicted 2-adic valuation pattern
$c_2(m)/\#H^1_{f,2}(m) \sim 2^{\{-3,-2,0,1\}}$ would be a striking
numerical consequence of the Bloch-Kato formula in the Steinberg-ramified case.}

%% ============================================================
\section{Why existing TNC machinery is blocked}
\label{sec:blocked}
%% ============================================================

We document the complete obstruction from published 2024-2026 papers.

\subsection{Büyükboduk-Neamti 2026}

Büyükboduk and Neamti \cite{BN26} prove a Tamagawa Number Conjecture
result via a Bertolini-Darmon-Prasanna style formula combined with
Beilinson-Kato elements. Their framework requires a list of hypotheses,
stated explicitly in §2.5.7 of \cite{BN26}:

\begin{quote}
\textbf{(Reg)} $k \geq 2$ and $h$ is a non-critical $p$-regular eigenform.\\
\textbf{(Eis)} $k = 1$ and $h$ is the $p$-stabilisation of a $p$-irregular Eisenstein series of weight $1$ and level $N_h$.\\
\textbf{(Disc)} $K \neq \Q(\sqrt{-1}),\, \Q(\sqrt{-3})$, and $\mathrm{disc}(K/\Q)$ is odd.\\
\textbf{(Heeg)} All rational primes dividing $N_h$ are split in $K$.\\
\textbf{(Sign)} The global root number of the base change of $h$ to $K$ equals $-1$.
\end{quote}

(Verbatim OCR of arXiv:2604.13854v1, §2.5.7, page~9; vision-verified by M69 on 2026-05-06.)

For $f = 4.5.b.a$:
\begin{itemize}
  \item \textbf{(Disc) first clause}: $K = \Q(i) = \Q(\sqrt{-1})$ is \emph{explicitly excluded by name}.
    The authors note that fields with additional automorphisms require separate treatment.
  \item \textbf{(Disc) second clause}: $\mathrm{disc}(K/\Q) = -4$ is even; hypothesis fails.
  \item \textbf{(Heeg)}: $2 \mid N = 4$, but $2$ ramifies in $K = \Q(i)$: $(2) = (1+i)^2$. Hypothesis fails.
  \item §1 \textbf{$p > 3$ split}: ``let $p > 3$ be a prime that splits in $K/\Q$''. For us $p = 2$: fails.
  \item §1 \textbf{$\mathrm{ord}_p(N_f) \leq 1$}: $\mathrm{ord}_2(4) = 2$: fails.
\end{itemize}

There are \emph{five} compounding hypothesis violations.
\cite{BN26} does not apply to $f = 4.5.b.a$.

\subsection{Other 2024-2025 TNC papers}

\begin{table}[h]
\centering
\renewcommand{\arraystretch}{1.3}
\begin{tabular}{lllll}
\hline
Paper & $p > 3$ split & weight even & adjoint only & $p \nmid N$ \\
\hline
Büyükboduk-Neamti \cite{BN26} & yes & no & no & yes \\
Castella 2024 \cite{Cas24} & yes (PLMS) & no & no & yes \\
BBL 2024 \cite{BBL24} & yes ($p \geq 5$) & no & no & yes \\
Sano 2025 \cite{San25} & yes (split) & yes (even) & no & yes \\
Yin 2024 \cite{Yin24} & weak (Eisenstein) & no & no & yes \\
Diamond-Flach-Guo \cite{DFG25} & weak & no & \emph{yes} & weak \\
\hline
\end{tabular}
\caption{Obstruction table for $f = 4.5.b.a$ with $p=2$ ramified, $k=5$ odd, $N=4$.
For each paper: ``yes'' in a column means the hypothesis is present in the paper
and fails for $f$. Diamond-Flach-Guo is blocked because it treats only the adjoint motive $\mathrm{Sym}^2 f$.}
\label{tab:obstructions}
\end{table}

\begin{itemize}
  \item \textbf{Castella \cite{Cas24}}: anticyclotomic Euler systems for supersingular elliptic curves, $p > 3$ split. Neither $p=2$ nor $p \mid N$ is treated.
  \item \textbf{Burungale-Büyükboduk-Lei \cite{BBL24}}: explicit abstract ``$p \geq 5$ supersingular''; does not cover $p=2$ ramified.
  \item \textbf{Sano \cite{San25}}: the main result requires $k = 2r$ even; $k=5$ is odd.
  \item \textbf{Yin \cite{Yin24}}: good Eisenstein framework requires $p \nmid N$; $2 \mid 4$ fails.
  \item \textbf{Diamond-Flach-Guo \cite{DFG25}}: treats only the adjoint motive $\mathrm{Sym}^2 f$, not $M(f)$ itself.
\end{itemize}

\textbf{Conclusion.} There is currently \emph{no published machinery} for the BK Tamagawa conjecture in the regime $(p=2\text{ ramified},\ k=5\text{ odd},\ K=\Q(i),\ N=4)$.

\subsection{What would be required}

Resolving Conjectures~R-2.1--R-2.3 requires a new framework combining:
\begin{enumerate}[label=(\arabic*)]
  \item \emph{Anticyclotomic Iwasawa Main Conjecture at $p=2$ ramified in $K$}:
    all published results (Hsieh \cite{Hsi14}, Chida-Hsieh \cite{CH18}, Arnold, Pollack-Weston)
    exclude $p = 2$ or $p$ ramified in $K$. The IMC is currently FALSIFIED for our regime.
  \item \emph{Existence of $\mathscr{L}_2^\pm(f)$} (Conjecture~M13.1.A in the ECI programme):
    the 2-adic anticyclotomic $L$-function must be constructed first (Fan-Wan \cite{FW23},
    Kriz \cite{Kriz21} extensions needed).
  \item \emph{Tate local cohomology at the Steinberg-ramified prime $p=2$}: TBD-R2-1.
\end{enumerate}

\tbdtag{5}{Reduce Conjecture~R-2.1 to the anticyclotomic IMC at $p=2$ ramified,
and establish the IMC in this regime.
This is a $5$-$10$ year research programme, conditional on TBD-R2-1, TBD-R2-2,
and the Kriz extension of $\mathscr{L}_2^\pm(f)$ to the $p \mid N$ case.}

%% ============================================================
\section{Compatibility with prior ECI modules}
\label{sec:compat}
%% ============================================================

We verify that Conjectures~R-2.1--R-2.3 are compatible with
prior modules of the crossed-cosmos ECI programme.

\subsection{M27 — Hodge conjecture (trivial)}

By Pohlmann's theorem \cite{Poh68}, all Hodge cycles on abelian varieties with
CM are absolute Hodge cycles, hence the Hodge conjecture is trivially
satisfied for $M(f)_B$. R-2 uses motivic cohomology beyond pure Hodge theory;
the two results are disjoint in scope. Verdict: \textbf{compatible, unchanged}.

\subsection{M31 — BSD subsumed (weight-2 only)}

The BSD conjecture in rank-$0$ form concerns weight-$2$ newforms
(elliptic curves). Module M31 established that the BSD contribution
of $f = 4.5.b.a$ is subsumed by M13.1 paper-2 (Hecke-Maass companion).
R-2 reopens the higher-weight BK Tamagawa conjecture as a \emph{standalone}
programme for $k = 5$, orthogonal to the BSD subclassification.
Verdict: \textbf{compatible, no conflict}.

\subsection{M39 — Beilinson regulator (non-critical $s=k=5$)}

Module M39 considers Beilinson's conjecture at the \emph{non-critical}
value $s = k = 5$, where the motivic cohomology regulator
$\langle r_{\mathcal{D}}(\xi_m^{KLZ}), \omega_f \rangle / \Omega_f$
appears. This is \emph{disjoint} from R-2, which concerns the critical
integers $m \in \{1,2,3,4\}$ where $L(f,m) \neq 0$ and the regulator
contribution vanishes (Beilinson's regulator is trivial at critical integers
for CM motives by a weight argument).
Verdict: \textbf{disjoint, compatible}.

\subsection{M52 — Damerell $6/5$ invariant}

Module M52 is the \emph{source} of the R-2 central data:
the $\Omega$-independent ratio $6/5$ and the full-$\Q$ Damerell ladder.
R-2 proposes a Bloch-Kato interpretation of what M52 computes.
Verdict: \textbf{M52 is prerequisite for R-2}.

\subsection{M55 — 4-layer uniqueness}

Module M55 identifies four structural layers uniquely characterizing $4.5.b.a$.
The Damerell-Tamagawa ratio conjecture (R-2.2) would provide a \emph{fifth layer}:
BK rationality of the $6/5$ invariant.
Verdict: \textbf{R-2 adds a 5th uniqueness layer to M55}.

\subsection{M57 — Adelic Katz conjecture}

Module M57 conjectures the existence of $\mathscr{L}_2^\pm(f) \in D(\Gamma, \Z_2)$
(Conjecture~M13.1.A). This is a \emph{necessary first step} for R-2.5.
Together, (M57 + R-2) would form an IMC pair at $p = 2$ ramified,
analogous to the Bloch-Kato + Iwasawa Main Conjecture pair for ordinary primes.
Verdict: \textbf{M57 is a necessary prerequisite for TBD-R2-5}.

%% ============================================================
\section{Open problems}
\label{sec:open}
%% ============================================================

We collect the five \textbf{[TBD: prove]} markers with difficulty estimates.

\begin{center}
\renewcommand{\arraystretch}{1.4}
\begin{tabular}{lll}
\hline
Tag & Description & Difficulty \\
\hline
TBD-R2-1 & $c_2(M(f)(m))$ at ramified Steinberg $p=2$ (Tate local $\varepsilon$-factors) & HIGH \\
TBD-R2-2 & $\#H^1_f(\Q, V_2(f)(m))$ at $m=1,2,3,4$ (no Heegner/Kolyvagin available) & HIGH \\
TBD-R2-3 & Factor $6/5 = c_2(1)/c_2(2) \times \mathrm{Selmer} \times \Gamma$-ratio & \textbf{VERY HIGH} \\
TBD-R2-4 & Deligne periods $\Omega^\pm_m$ via Chowla-Selberg + lemniscate & MEDIUM \\
TBD-R2-5 & Reduce to M57 + anticyclotomic IMC at $p=2$ ramified & HIGH (IMC blocked) \\
\hline
\end{tabular}
\end{center}

\medskip
\noindent\textbf{TBD-R2-3 is the primary open problem}. Its resolution would constitute
a non-trivial verification of the Bloch-Kato formula~\eqref{eq:R21} for a CM motive
at a Steinberg-edge ramified prime. The factorization \eqref{eq:R22} is purely arithmetic,
requiring no period choices, and is in principle accessible to direct Selmer computation
and local Tate cohomology once TBD-R2-1 and TBD-R2-2 are resolved.

\medskip
\noindent\textbf{Potential collaborators}: Ivan Kriz (MIT) for TBD-R2-1
(Hodge-filtration extension, $p \mid N$ case); Büyükboduk (UC Dublin) for
BDP extension to $p=2$ ramified; Antonio Lei (Ottawa) for the anticyclotomic
IMC framework (TBD-R2-5); Castella (UCSB) for anticyclotomic Selmer at $p=2$.

%% ============================================================
\section*{Declarations}
%% ============================================================

\noindent\textbf{Conflict of interest.} None.

\noindent\textbf{Data availability.} All numerical data in this note
is reproducible from the PARI/GP and SageMath scripts available from
the corresponding author. LMFDB data is publicly accessible at
\url{https://www.lmfdb.org/ModularForm/GL2/Q/holomorphic/4/5/b/a/}.

\noindent\textbf{Honesty bar.} This note contains \emph{conjectures},
not theorems. The five [TBD: prove] markers are explicit.
The honest probability estimate of $\sim 10$-$15\%$ for a formal
Bloch-Kato Tamagawa contribution is stated in the abstract
and derived from the complete obstruction documented in Section~\ref{sec:blocked}.

%% ============================================================
\begin{thebibliography}{99}
%% ============================================================

\bibitem{BK90}
S.~Bloch and K.~Kato,
\emph{$L$-functions and Tamagawa numbers of motives},
in: \emph{The Grothendieck Festschrift, Vol. I},
Progr. Math. \textbf{86}, Birkhäuser, Boston, 1990, pp.~333--400.
DOI:~\href{https://doi.org/10.1007/978-0-8176-4574-8_9}{10.1007/978-0-8176-4574-8\_9}.

\bibitem{FPR94}
J.-M.~Fontaine and B.~Perrin-Riou,
\emph{Autour des conjectures de Bloch et Kato: cohomologie galoisienne
et valeurs de fonctions $L$},
in: \emph{Motives} (Seattle, WA, 1991),
Proc. Sympos. Pure Math. \textbf{55}, Part~1,
Amer. Math. Soc., Providence, RI, 1994, pp.~599--706.

\bibitem{Del79}
P.~Deligne,
\emph{Valeurs de fonctions $L$ et périodes d'intégrales},
in: \emph{Automorphic Forms, Representations, and $L$-Functions},
Proc. Sympos. Pure Math. \textbf{33}, Part~2,
Amer. Math. Soc., 1979, pp.~313--346.

\bibitem{Del71}
P.~Deligne,
\emph{Formes modulaires et représentations $\ell$-adiques},
in: \emph{Séminaire Bourbaki}, Exp. No.~355,
Lecture Notes in Math. \textbf{179}, Springer, 1971, pp.~139--172.

\bibitem{Sch90}
A.J.~Scholl,
\emph{Motives for modular forms},
Invent. Math. \textbf{100} (1990), 419--430.
DOI:~\href{https://doi.org/10.1007/BF01231196}{10.1007/BF01231196}.

\bibitem{CS67}
S.~Chowla and A.~Selberg,
\emph{On Epstein's zeta-function},
J. Reine Angew. Math. \textbf{227} (1967), 86--110.
DOI:~\href{https://doi.org/10.1515/crll.1967.227.86}{10.1515/crll.1967.227.86}.

\bibitem{Dam70}
R.M.~Damerell,
\emph{$L$-functions of elliptic curves with complex multiplication, I},
Acta Arith. \textbf{17} (1970), 287--301.
DOI:~\href{https://doi.org/10.4064/aa-17-3-287-301}{10.4064/aa-17-3-287-301}.

\bibitem{Dam71}
R.M.~Damerell,
\emph{$L$-functions of elliptic curves with complex multiplication, II},
Acta Arith. \textbf{19} (1971), 311--317.
DOI:~\href{https://doi.org/10.4064/aa-19-3-311-317}{10.4064/aa-19-3-311-317}.

\bibitem{Shim76}
G.~Shimura,
\emph{The special values of the zeta functions associated with cusp forms},
Comm. Pure Appl. Math. \textbf{29} (1976), 783--804.
DOI:~\href{https://doi.org/10.1002/cpa.3160290618}{10.1002/cpa.3160290618}.

\bibitem{Katz78}
N.M.~Katz,
\emph{$p$-adic $L$-functions for CM fields},
Invent. Math. \textbf{49} (1978), 199--297.
DOI:~\href{https://doi.org/10.1007/BF01390187}{10.1007/BF01390187}.

\bibitem{CO77}
H.~Cohen and J.~Oesterlé,
\emph{Dimensions des espaces de formes modulaires},
in: \emph{Modular Functions of One Variable VI},
Lecture Notes in Math. \textbf{627}, Springer, 1977, pp.~69--78.

\bibitem{Rub91}
K.~Rubin,
\emph{The ``main conjectures'' of Iwasawa theory for imaginary quadratic fields},
Invent. Math. \textbf{103} (1991), 25--68.
DOI:~\href{https://doi.org/10.1007/BF01239508}{10.1007/BF01239508}.

\bibitem{Poh68}
H.~Pohlmann,
\emph{Algebraic cycles on abelian varieties of complex multiplication type},
Ann. of Math. \textbf{88} (1968), 161--180.
DOI:~\href{https://doi.org/10.2307/1970570}{10.2307/1970570}.

\bibitem{BN26}
K.~Büyükboduk and A.~Neamti,
\emph{Bloch-Kato Selmer groups and the Tamagawa Number Conjecture for
modular forms via $p$-adic Gross-Zagier formulae},
arXiv:2604.13854 (2026).

\bibitem{Cas24}
F.~Castella,
\emph{Anticyclotomic Euler systems for unitary Shimura varieties},
Proc. London Math. Soc. (2025), arXiv:2407.11891.

\bibitem{BBL24}
A.~Burungale, K.~Büyükboduk, and A.~Lei,
\emph{Perrin-Riou's conjecture for supersingular elliptic curves},
arXiv:2310.06813 (2024); to appear in Algebra Number Theory.

\bibitem{San25}
K.~Sano,
\emph{Bloch-Kato conjecture for symmetric square motives},
arXiv:2510.01601 (2025).

\bibitem{Yin24}
C.~Yin,
\emph{Iwasawa theory for good Eisenstein primes},
arXiv:2410.24193 (2024).

\bibitem{DFG25}
F.~Diamond, M.~Flach, and L.~Guo,
\emph{Adjoint Selmer groups and Bloch-Kato conjecture for $\mathrm{GL}_2$},
arXiv:2512.02348 (2025).

\bibitem{Kriz21}
I.~Kriz,
\emph{Supersingular $p$-adic $L$-functions, Maass-Shimura Operators
and Waldspurger Formulas},
Princeton/AMS Monographs, vol.~212, 2021.
DOI:~\href{https://doi.org/10.1515/9780691225739}{10.1515/9780691225739}.

\bibitem{Hsi14}
M.-L.~Hsieh,
\emph{Special values of anticyclotomic Rankin-Selberg $L$-functions},
Doc. Math. \textbf{19} (2014), 709--767.

\bibitem{CH18}
M.~Chida and M.-L.~Hsieh,
\emph{On the anticyclotomic Iwasawa main conjecture for modular forms},
Compos. Math. \textbf{151} (2018), no.~5, 863--900.
DOI:~\href{https://doi.org/10.1112/S0010437X14007787}{10.1112/S0010437X14007787}.

\bibitem{FW23}
X.~Fan and L.~Wan,
\emph{$p$-adic $L$-functions for $\mathrm{GL}_2$ at Steinberg primes},
arXiv:2304.09806 (2023).

\bibitem{LMFDB}
The LMFDB Collaboration,
\emph{The $L$-functions and Modular Forms Database},
\url{https://www.lmfdb.org} (2024).
Form 4.5.b.a: \url{https://www.lmfdb.org/ModularForm/GL2/Q/holomorphic/4/5/b/a/}.

\bibitem{PARI}
The PARI Group,
\emph{PARI/GP version 2.15.x},
Bordeaux, 2024, \url{https://pari.math.u-bordeaux.fr/}.

\bibitem{Sage}
W.A.~Stein et al.,
\emph{Sage Mathematics Software (Version 10.x)},
The Sage Development Team, 2024, \url{https://www.sagemath.org}.

\end{thebibliography}

\end{document}
